\chapter{Chapter05 — Linear Threshold Influence}
%%% generated by the blueprint pipeline from FABL.Chapter05.LinearThresholdInfluence
\section{Overview}

Declarations of the Lean module \texttt{FABL.Chapter05.LinearThresholdInfluence}.

\section{Declarations}

\begin{theorem}[Threshold signs of a sum and difference]
\lean{FABL.thresholdSign_add_ne_sub_iff}
\label{FABL.thresholdSign_add_ne_sub_iff}
\leanok
\uses{FABL.thresholdSign, FABL.thresholdSign_of_neg, FABL.thresholdSign_of_nonneg}
\difficulty{5}
For all real numbers $r$ and $a$, the threshold signs of $r+a$ and $r-a$ differ if and
only if $-|a|\le r<|a|$.
\end{theorem}

\begin{definition}[Affine linear remainder]
\lean{FABL.affineLinearRemainder}
\label{FABL.affineLinearRemainder}
\leanok
\uses{FABL.signValue}
Given $a_0\in\mathbb R$, coefficients $a_k\in\mathbb R$ for $k$ in a set of $n$
coordinates, a coordinate $i$, and a sign vector $x\in\{-1,1\}^n$, the affine linear
remainder at $i$ is $a_0+\sum_{k\ne i}a_k x_k$.
\end{definition}

\begin{definition}[Two-coordinate affine remainder]
\lean{FABL.affineLinearRemainderTwo}
\label{FABL.affineLinearRemainderTwo}
\leanok
\uses{FABL.signValue}
Given $a_0\in\mathbb R$, coefficients $a_k\in\mathbb R$ for $k$ in a set of $n$
coordinates, coordinates $i,j$, and a sign vector $x\in\{-1,1\}^n$, the two-coordinate
affine remainder is $a_0+\sum_{k\ne i,\ k\ne j}a_k x_k$.
\end{definition}

\begin{theorem}[Affine form after setting one coordinate]
\lean{FABL.affineLinearForm_setCoordinate}
\label{FABL.affineLinearForm_setCoordinate}
\leanok
\uses{FABL.Sign, FABL.affineLinearForm, FABL.affineLinearRemainder, FABL.linearForm, FABL.setCoordinate, FABL.setCoordinate_apply_of_ne, FABL.setCoordinate_apply_self, FABL.signValue}
\difficulty{4}
Let $a_0\in\mathbb{R}$, let $a\in\mathbb{R}^n$, let $i$ be a coordinate, let
$x\in\{-1,1\}^n$, and let $s\in\{-1,1\}$. The affine linear form evaluated at the vector
obtained from $x$ by setting coordinate $i$ to $s$ equals $a_0+\sum_{k\ne
i}a_kx_k+a_is$.
\end{theorem}

\begin{theorem}[Decomposition of an affine linear remainder]
\lean{FABL.affineLinearRemainder_eq_two_add}
\label{FABL.affineLinearRemainder_eq_two_add}
\leanok
\uses{FABL.affineLinearRemainder, FABL.affineLinearRemainderTwo, FABL.signValue}
\difficulty{3}
Let $a_0\in\mathbb{R}$, let $a\in\mathbb{R}^n$, let $i,j$ be distinct coordinates, and
let $x\in\{-1,1\}^n$. The affine linear remainder obtained by omitting coordinate $j$
equals the remainder obtained by omitting both $i$ and $j$, plus the term $a_ix_i$.
\end{theorem}

\begin{theorem}[Coordinate swap identity for an affine remainder]
\lean{FABL.affineLinearRemainder_permute_swap_eq}
\label{FABL.affineLinearRemainder_permute_swap_eq}
\leanok
\uses{FABL.affineLinearRemainder, FABL.affineLinearRemainderTwo, FABL.permuteInput, FABL.signValue}
\difficulty{4}
Let $n$ be a natural number, let $a_0\in\mathbb R$ and $a_1,\ldots,a_n\in\mathbb R$, and
let $i,j$ be distinct coordinates in $\{1,\ldots,n\}$. For a sign vector
$x\in\{-1,1\}^n$, write $x^{(ij)}_k=x_{(ij)(k)}$, where $(ij)$ swaps $i$ and $j$, and
interpret each sign as its real value. The sum $a_0+\sum_{k\ne i}a_k x^{(ij)}_k$ equals
$a_0+\sum_{k\notin\{i,j\}}a_k x_k+a_jx_i$.
\end{theorem}

\begin{theorem}[Threshold sign multiplication identity]
\lean{FABL.mul_signValue_thresholdSign_mul}
\label{FABL.mul_signValue_thresholdSign_mul}
\leanok
\uses{FABL.Sign, FABL.signValue, FABL.thresholdSign, FABL.thresholdSign_of_neg, FABL.thresholdSign_of_nonneg}
\difficulty{4}
Let $a\in\mathbb R$ and let $s\in\{ -1,1\}$. Define the sign of a real number $u$ to be
$1$ when $u\ge0$ and $-1$ when $u<0$. Then $a\,\operatorname{sign}(a)s=|a|s$.
\end{theorem}

\begin{theorem}[Pivotality of an affine threshold]
\lean{FABL.isPivotal_linearThreshold_iff}
\label{FABL.isPivotal_linearThreshold_iff}
\leanok
\uses{FABL.BooleanFunction, FABL.IsPivotal, FABL.affineLinearForm, FABL.affineLinearForm_setCoordinate, FABL.affineLinearRemainder, FABL.isPivotal_iff_setCoordinate_ne, FABL.thresholdSign, FABL.thresholdSign_add_ne_sub_iff}
\difficulty{3}
Let $n\in\mathbb N$, let $f:\{\pm1\}^n\to\{\pm1\}$, and suppose that $f(x)$ is the sign
of $a_0+\sum_{k=1}^n a_k x_k$ for every $x$, where the sign is $+1$ for a nonnegative
argument and $-1$ otherwise. For a coordinate $i$ and input $x$, write
$R_i(x)=a_0+\sum_{k\ne i}a_kx_k$. Then $f$ is pivotal at $x$ in coordinate $i$ if and
only if $-|a_i|\le R_i(x)<|a_i|$.
\end{theorem}

\begin{theorem}[Boolean influence le of nonnegative weight le]
\lean{FABL.booleanInfluence_le_of_nonnegative_weight_le}
\statementsource{Analysis-of-Boolean-Functions-by-Ryan-ODonnell}{pdf-ffbe9c7b3306-p134-b005}
\label{FABL.booleanInfluence_le_of_nonnegative_weight_le}
\leanok
\uses{FABL.BooleanFunction, FABL.IsPivotal, FABL.affineLinearForm, FABL.affineLinearRemainder_eq_two_add, FABL.affineLinearRemainder_permute_swap_eq, FABL.booleanInfluence, FABL.isPivotal_linearThreshold_iff, FABL.permuteInput, FABL.permuteInputEquiv, FABL.signValue, FABL.thresholdSign, FABL.uniformProbability}
\difficulty{6}
% dataset-pipeline: stub -- prose pending build_informal_statements.py
Statement pending regeneration.
\end{theorem}

\begin{theorem}[Influence comparison by threshold weights]
\lean{FABL.influence_le_of_abs_linearThresholdWeight_le}
\proofsource{Analysis-of-Boolean-Functions-by-Ryan-ODonnell}{
  pdf-ffbe9c7b3306-p134-b005,
  pdf-ffbe9c7b3306-p132-b002
}
\label{FABL.influence_le_of_abs_linearThresholdWeight_le}
\leanok
\uses{FABL.BooleanFunction, FABL.BooleanFunction.toReal, FABL.affineLinearForm, FABL.booleanInfluence, FABL.booleanInfluence_comp_reorientInput, FABL.booleanInfluence_eq_influence_toReal, FABL.booleanInfluence_le_of_nonnegative_weight_le, FABL.influence, FABL.linearForm, FABL.mul_signValue_thresholdSign_mul, FABL.thresholdSign}
\difficulty{4}
Let $f:\{-1,1\}^n\to\{-1,1\}$ be represented as
$f(x)=\operatorname{sgn}(a_0+\sum_{k=1}^n a_kx_k)$, with the threshold sign equal to $1$
when its argument is nonnegative and $-1$ otherwise. For any coordinates $i,j$, if
$|a_i|\leq |a_j|$, then $\operatorname{Inf}_i(f)\leq\operatorname{Inf}_j(f)$.
\end{theorem}

\begin{theorem}[Influence lower bound for balanced linear thresholds]
\lean{FABL.exists_one_div_sqrt_two_mul_dimension_le_influence_of_balanced_linearThreshold}
\label{FABL.exists_one_div_sqrt_two_mul_dimension_le_influence_of_balanced_linearThreshold}
\leanok
\uses{FABL.BooleanFunction, FABL.BooleanFunction.toReal, FABL.IsBalanced, FABL.IsLinearThreshold, FABL.IsUnate, FABL.abs_fourierCoeff_singleton_eq_influence_iff_isUnateInCoordinate, FABL.fourierCoeff, FABL.fourierWeightAtMost_one_eq_empty_add_sum_singleton, FABL.influence, FABL.influence_nonneg, FABL.isBalanced_iff_fourierCoeff_empty_eq_zero, FABL.isUnate_of_isLinearThreshold, FABL.one_half_le_fourierWeightAtMost_one_of_isLinearThreshold}
\difficulty{5}
Let $f:\{−1,1\}^n\to\{−1,1\}$ be balanced and representable as a linear threshold
function, with $n\ge1$. Then some coordinate $i$ has influence at least $1/\sqrt{2n}$.
\end{theorem}
